% ===================================================================
%  Biologie: exponentielles und logistisches Populationswachstum
%  aus denselben drei Parametern – BERECHNET, mit reinem TikZ.
%  Kompilieren: pdflatex vorlage_wachstum.tex
% ===================================================================
\documentclass[border=2mm]{standalone}
\usepackage[T1]{fontenc}
\usepackage{lmodern}
\usepackage[ngerman]{babel}
\usepackage{amsmath}
\usepackage{tikz}

\definecolor{akzentblau}{HTML}{2563EB}
\definecolor{akzentrot}{HTML}{DC2626}

% ---- Parameter -----------------------------------------------------
\def\Nnull{10}     % Startpopulation N_0
\def\rate{0.5}     % Wachstumsrate pro Tag
\def\K{500}        % Kapazität des Lebensraums
% --------------------------------------------------------------------

\begin{document}
\begin{tikzpicture}[x=0.45cm, y=0.01cm, font=\small]
  \pgfmathsetmacro\ymax{1.3*\K}
  % Ende der Exponentialkurve: dort, wo sie den oberen Rand erreicht
  \pgfmathsetmacro\tend{min(20, ln(\ymax/\Nnull)/\rate)}
  % Gitter und Achsen
  \draw[dashed, black!20] (0,0) grid[xstep=5, ystep=100] (20,\ymax);
  \draw[->] (0,0) -- (21,0);
  \node[below=18pt] at (10,0) {Zeit $t$ in Tagen};
  \draw[->] (0,0) -- (0,\ymax+60) node[above] {Individuen $N$};
  \foreach \x in {0,5,...,20} \node[below] at (\x,0) {\x};
  \foreach \y in {100,200,...,600} \node[left] at (0,\y) {\y};
  % Kapazität, exponentiell, logistisch
  \draw[black!60, dotted, thick] (0,\K) -- (20,\K) node[above left, black!70] {Kapazität $K=\K$};
  \draw[akzentrot, very thick, dashed] plot[domain=0:\tend, samples=150] (\x,{\Nnull*exp(\rate*\x)});
  \draw[akzentblau, very thick] plot[domain=0:20, samples=200] (\x,{\K/(1 + (\K/\Nnull - 1)*exp(-\rate*\x))});
  % Legende
  \draw[akzentrot, very thick, dashed] (13.5,180) -- (15,180) node[right, text=black] {exponentiell};
  \draw[akzentblau, very thick] (13.5,110) -- (15,110) node[right, text=black] {logistisch};
\end{tikzpicture}
\end{document}
